Partindo da relação de potência:
\[
y=mx^b,\quad m,b\in\mathbb{R},\; x>0.
\]

\textbf{Passo 1 — aplicar logaritmo em ambos os lados:}
\[
\log y = \log(mx^b)=\log m+\log x^b = \log m + b\log x.
\]

\textbf{Passo 2 — substituição de variáveis:} define-se
\[
Y=\log y,\quad X=\log x,\quad c=\log m.
\]
A equação transforma-se em:
\[
Y=bX+c,
\]
que é a equação de uma \emph{reta} no plano $(X,Y)$, com inclinação $b$ e intercepto $c=\log m$.

\textbf{Recuperação dos parâmetros originais:}
\begin{itemize}
  \item A inclinação da reta lida no gráfico log-log é diretamente o expoente $b$.
  \item O intercepto $c$ dá $m=10^c$ (para logaritmo de base 10).
\end{itemize}

\textbf{Exemplo numérico:} para $y=3x^2$ tem-se $m=3$, $b=2$.
\[
X=\log x,\quad Y=\log y = \log 3 + 2\log x = 0{,}477 + 2X.
\]
A reta log-log tem inclinação $2$ e intercepto $0{,}477$, confirmando $m=10^{0{,}477}\approx3$. $\checkmark$

